\documentclass[12pt, a4paper]{article}
\usepackage[utf8]{inputenc}
\usepackage{amsmath, amssymb, amsthm}
\usepackage{graphicx}
\usepackage{hyperref}
\usepackage{geometry}
\geometry{margin=1in}

\usepackage{tikz}
\usetikzlibrary{positioning, arrows.meta}

\newtheorem{definition}{Definition}[section]
\newtheorem{assumption}{Assumption}[section]
\newtheorem{proposition}{Proposition}[section]

\title{Time: The Undefined Parameter \\
\large Toward a Relational Construction of Temporal Order}
\author{
  Jerry Z. Xie \\
  \small The Silicon Valley Research, Inc. \\
  \small 2903 Bunker Hill Lane, Santa Clara, CA 95054, USA \\
  \small \texttt{jerry.z.xie@svrinc.io}
}
\date{}

\begin{document}

\maketitle

\begin{abstract}
Physics has used a temporal parameter for over three centuries. Newton used it. Thermodynamics used it. Electrodynamics used it. Einstein used it. Quantum mechanics uses it. Within their respective formalisms, this parameter has consistently been treated as primitive---assumed, transformed, parametrized, but not operationally defined.

This paper does not attempt to define time directly. Instead, it asks a constructive question: given a set of physical events and their relations, can we construct an ordering that plays the functional role of time---without presupposing a time parameter?

We formulate a relational structure in which admissible causal relations define a partial order over distinguishable events. Persistent records---defined as recoverable mutual information between degrees of freedom---provide a mechanism by which a subsystem can internally represent portions of that order. This yields a layered distinction between causal order, temporal order (as internally reconstructed), and numerical clock time (as a monotonic parametrization of that reconstruction).

We distinguish three questions that are often conflated: what orders events, how is that order represented, and how is it parametrized. The framework is consistent with relativity and with causal set theory, to which we acknowledge its debt. We do not claim that time ``is'' this structure. We claim that this structure is sufficient to construct a temporal ordering and its numerical representation, thereby recovering an important functional role normally played by time in physical description.
\end{abstract}

\section{Introduction: A Question I Was Afraid to Ask}

Fifty years ago, I stood at the lectern of a university physics department, teaching undergraduate quantum mechanics.

It was a lecture on the Schr\"odinger equation. I wrote on the blackboard:
\[
i\hbar \frac{\partial \psi}{\partial t} = \hat{H}\psi
\]
I stopped. I looked at that \( t \)---time. A voice in my head said: \textit{What is this thing, exactly?}

But the lecture had to continue. I did not ask.

In the decades that followed, I did not ask anyone. Not because I was not curious. Because the question was too foundational. Time appears throughout the standard equations of physics, yet its status is often left primitive within the formalism. And I---a person standing at a lectern teaching quantum mechanics---felt I \textit{should} know. I did not. And I lacked the courage to say, ``I do not know.''

Today, fifty years later, I still cannot say I know. But at least I am no longer afraid to ask.

The question of time has been asked philosophically for millennia---by Aristotle, Leibniz, Mach, and many others. Within physics, however, the temporal parameter has generally been treated as primitive within working formalisms. This paper does not ask the philosophical question anew. It asks a constructive one: \textbf{what physical structure is sufficient to recover the functional role of time, without taking a time parameter as given?}

\section{Newton: Time as Stage}

Newton wrote: \textit{``Absolute, true, and mathematical time, of itself and from its own nature, flows equably without relation to anything external.''}

Time, in this picture, is a \textbf{stage}. Matter moves; time flows independently. The stage is not affected by what happens upon it. This is the first major role time plays in physics: \textbf{the default container}.

A container, however, is not a definition. Newton tells us what time \textit{does}---it flows equably. He does not tell us what it \textit{is}.

\section{Thermodynamics: Direction, But No Definition}

Thermodynamics introduces the concept of entropy. The second law introduces an empirical arrow of time: under appropriate macroscopic conditions, entropy tends to increase toward what is conventionally called the future.

But thermodynamics does not define time. It identifies a \textit{property} of time---that it has a direction. If one asks thermodynamics ``What is time?'' it answers ``It has a direction.'' This describes a behavior, not an essence.

\textbf{The contribution of thermodynamics:} time has an arrow---but an arrow is not a definition.

\section{Electrodynamics: Time as Field Evolution Parameter}

Maxwell's equations unified electricity and magnetism. Classical electrodynamics was initially formulated within the Newtonian framework of absolute time. In this framework, time is the parameter with respect to which electromagnetic fields evolve. However, the invariant propagation speed appearing in Maxwell's equations was incompatible with Galilean relativity and ultimately contributed to the development of special relativity.

The equations assume a universal time coordinate, but embedded within them is the speed of light---a constant that would later be used to break the assumption of absolute time.

\textbf{The contribution of electrodynamics:} time became the background parameter for field evolution---still assumed, still undefined.

\section{Einstein: Time as Coordinate}

Einstein showed that time cannot be treated as absolute. It merges with space into spacetime.

Time becomes a \textbf{coordinate}---the fourth component of a four-dimensional manifold. Its measured intervals depend on the spacetime geometry and on the observer's worldline. Relativity provides a precise geometric and operational definition of proper time along timelike worldlines. But this does not by itself answer the more foundational question of whether the temporal parameter is fundamental, or why the underlying spacetime/causal structure exists.

\textbf{The contribution of Einstein:} time is not a stage; it is part of the geometry. Physics can define particular manifestations of time without explaining why temporal order exists fundamentally.

\section{Quantum Mechanics: Time as Parameter, and Nothing More}

In standard nonrelativistic quantum mechanics, time is not represented by a self-adjoint operator in the same way as position or momentum. This is related to the well-known ``problem of time'' that becomes acute in attempts to quantize gravity.

The Schr\"odinger equation contains a \( t \). But that \( t \) is a number, not a measured quantity in the same sense as position or momentum. Time therefore occupies a distinctive status in the standard formalism: it serves as an evolution parameter, while physical durations are operationally compared using clocks.

\textbf{The contribution of quantum mechanics:} time occupies a special status---it is a parameter, not an observable in the usual sense.

\section{The Common Thread: All Roles, No Definition}

Each framework gives time a new role. None gives it a definition. This pattern suggests something structural: time is not an object within physics, but a condition under which physics operates.

\section{A Constructive Framework for Temporal Order}

We now ask a different question: \textbf{Given a set of physical events and their relations, can we construct an ordering that functions as time---without presupposing a time parameter?}

The framework is relational. It does not claim that time ``is'' these relations. It claims that these relations are \textit{sufficient} to construct a temporal ordering and its numerical representation.

\subsection{States and Events}

Let \( \mathcal{S} \) be a set of possible physical states. An \textbf{event} \( E \) is a transition:
\[
E: S_i \rightarrow S_j, \quad S_i \neq S_j
\]
We require that \( S_i \) and \( S_j \) are \textbf{distinguishable}: there exists an observable \( O \) such that \( O(S_i) \neq O(S_j) \).

\textbf{Definition:} The set of all events is \( \mathcal{E} \).

\subsection{Causal Order}

Let \( E_A, E_B \in \mathcal{E} \). We write:
\[
E_A \prec_C E_B
\]
to mean that \( E_B \) is physically dependent on \( E_A \)---there exists an admissible physical influence from \( A \) to \( B \), such that \( B \) cannot be fully described without reference to \( A \).

\textbf{Assumption (Causal Consistency):} The relation \( \prec_C \) is a strict partial order on \( \mathcal{E} \):

1. \textbf{Irreflexive:} \( E_A \nprec_C E_A \) for all \( E_A \).
2. \textbf{Transitive:} If \( E_A \prec_C E_B \) and \( E_B \prec_C E_C \), then \( E_A \prec_C E_C \).

Irreflexivity and transitivity together imply asymmetry. This is not derived from distinguishability. It is posited as a physical consistency condition---one consistent with relativity's light-cone structure and with the causal structure of classical and quantum field theories.

This causal order is the first layer. It is \textit{assumed}, not derived.

\subsection{Persistent Records}

Let \( E_A \) and \( E_B \) be causally related events, with \( E_A \prec_C E_B \). Let \( X_A \) represent the physical degrees of freedom associated with event \( E_A \). Let \( R_B \) represent the degrees of freedom in the downstream subsystem that encode information about \( E_A \).

\textbf{Definition (Persistent Record):} A persistent record exists if:
\[
I(X_A; R_B) > I_{\min}
\]
where \( I(X_A; R_B) \) is the mutual information between the degrees of freedom of the earlier event and the downstream encoding, and \( I_{\min} \) is an operational recoverability threshold.

Here \( I_{\min} \) is not assumed to be a universal constant; it represents an operational recoverability threshold associated with the relevant subsystem and measurement procedure. ``Persistent'' means that the correlation remains recoverable over a physically relevant interval or under a specified class of perturbations. It does not imply absolute permanence.

\textbf{Note:} Mutual information is symmetric: \( I(X_A; R_B) = I(R_B; X_A) \). It therefore cannot by itself specify the direction of the record. The direction is supplied by the causal relation \( \prec_C \). The information correlation provides the physical substrate of the record; the causal order provides its orientation.

\textbf{Interpretation:} A persistent record is a correlation between degrees of freedom that preserves information about an earlier event. Records do not create causal order. They provide the physical mechanism by which a subsystem can internally represent that order.

\subsection{Accessible Events and Temporal Order}

Causal order exists independently of observation. However, a physical subsystem may or may not have access to the full causal structure.

An \textbf{observer} \( \mathcal{O} \) is a physical subsystem that:
1. Has internal states,
2. Interacts with other subsystems,
3. Stores persistent records of those interactions.

Define the observer's \textbf{accessible event set} as:
\[
\mathcal{E}_{\mathcal{O}}
=
\left\{
E\in\mathcal{E}:
\exists R_{\mathcal O}(E)
\text{ such that }
I(X_E;R_{\mathcal O}(E))>I_{\min}
\right\}
\]
where \( R_{\mathcal O}(E) \) is the record of \( E \) stored by \( \mathcal O \), and \( X_E \) represents the degrees of freedom associated with \( E \).

In words: an event is accessible to the observer if and only if the observer has a recoverable persistent record of it.

\textbf{Definition (Temporal Order for \( \mathcal{O} \)):} The observer's temporal order \( \prec_{T,\mathcal{O}} \) is the restriction of the causal partial order \( \prec_C \) to accessible events:
\[
E_A \prec_{T,\mathcal{O}} E_B
\iff
E_A \prec_C E_B
\quad\text{and}\quad
E_A, E_B \in \mathcal{E}_{\mathcal{O}}.
\]

\textbf{Key distinction:}
- Causal order \( \prec_C \) exists regardless of observers.
- Temporal order \( \prec_{T,\mathcal{O}} \) is the portion of causal order that an observer can internally reconstruct through records.
- The observer does \textit{not} create causal order. The observer reconstructs a part of it.

We use the term \textit{temporal order} here in a functional sense: it is an ordering that can support the temporal descriptions used by a physical subsystem. We do not assume that causal order and experienced or operational time are identical in every physical theory.

\subsection{Clock Time: Parametrizing Order}

A physical clock does not typically parametrize all accessible events. It follows a specific worldline or history.

Let \( \mathcal H_{\mathcal O} \subseteq \mathcal E_{\mathcal O} \) be an ordered history along a clock's worldline.

A \textbf{clock} is a physical process that provides an order-preserving numerical parametrization:
\[
\tau: \mathcal H_{\mathcal O} \rightarrow \mathbb{R}
\]
such that for any two events in \( \mathcal H_{\mathcal O} \):
\[
E_A \prec_{T,\mathcal{O}} E_B \Rightarrow \tau(E_A) < \tau(E_B).
\]

\textbf{Definition:} This mapping \( \tau \) is a \textbf{numerical representation} of temporal order along the clock's history.

\textbf{The representation is not unique.} Any strictly increasing function \( f \) yields a different clock:
\[
\tau'(E) = f(\tau(E)), \qquad f \text{ strictly increasing}.
\]

This non-uniqueness means that the numerical time parameter carries more information than the underlying order---it includes choices of scale, origin, and parametrization. The underlying physical structure is the order, not the numbers.

We do not assume a single clock can assign an order-preserving numerical label to all accessible events simultaneously.

\subsection{The Arrow of Time}

Within this framework, the arrow of time is defined as follows:

\textbf{Definition:} The direction of time is the direction in which stable persistent records accumulate.

If \( E_A \prec_C E_B \) and \( E_B \) contains records of \( E_A \), then records form in the direction of the causal order. This direction is intrinsic to the record-forming process.

\textbf{Relationship to thermodynamics:} In macroscopic systems, the direction of record formation corresponds to the direction of entropy increase. This is a strong empirical correspondence, but the precise relationship between record formation and entropy production remains an open problem. We do not claim to have reduced the thermodynamic arrow to records---only to have identified a parallel structural arrow that operates at the level of information preservation.

\subsection{Consistency with Relativity}

In special and general relativity, the causal structure is determined by light cones. Events with a causal connection have an invariant order; events without such a connection (spacelike separated) do not.

In our framework:
- For causally connected events, \( E_A \prec_C E_B \) is invariant---it does not depend on the observer.
- For spacelike separated events, \( E_A \parallel E_B \), there is no invariant causal order. Different observers may assign different \( \tau \) values, and there is no physically meaningful ``before'' or ``after.''

This is precisely the content of relativity's ``relativity of simultaneity.'' Our framework is consistent with it, but does not attempt to derive it from records. Relativity provides the light-cone structure; the framework provides a relational substrate on which that structure can be represented and parametrized.

\subsection{Relation to Causal Set Theory}

This framework is closely related to causal set theory, which treats spacetime as a locally finite partial order of events whose order represents causality. In that approach, the partial order is taken as the fundamental structure from which continuum spacetime is recovered as an approximation.

Our framework shares this starting point. Its distinct contributions are:
1. The explicit introduction of \textbf{persistent records} (recoverable mutual information between degrees of freedom) as the mechanism by which an observer internally represents causal order.
2. The \textbf{layered distinction} between causal order, temporal order, and clock time.
3. The \textbf{thesis} that the numerical time parameter need not be fundamental---it can be treated as a representation of an underlying relational structure.

\section{Epilogue: A Boundary We Share with Newton}

Newton's natural philosophy ultimately rested on a theological framework that extended beyond the mechanics itself. When Newton reached the limit of what physics could explain, he invoked God. Not as a theological claim, but as a placeholder---something that accounted for what physics itself could not account for.

We now stand at a similar boundary. We have equations that use time. We have no equation that defines it. We have measurements that rely on time. We have no measurement that reveals what it is. We have theories that depend on time. We have no theory that explains why time exists at all.

Time occupies in modern physics the same position that ``God'' occupied for Newton: the unexamined ground upon which everything else is built.

This is not a religious claim. It is a structural observation. It identifies the position that time holds in our theoretical framework: not an object within physics, but a condition under which physics operates.

We do not need to invoke God. But we do need to acknowledge the boundary---and to stop pretending that we have already crossed it.

\section{Conclusion}

This paper has not answered ``what is time?'' It has asked a different, more constructive question: \textit{can we construct a structure sufficient to play the role of time, without presupposing a time parameter?}

We have proposed such a structure. It consists of:
1. \textbf{Causal order}, defined as a strict partial order over events, assumed as a physical consistency condition.
2. \textbf{Persistent records}, defined as recoverable mutual information between degrees of freedom, providing the mechanism for internal reconstruction.
3. \textbf{Temporal order}, the restriction of causal order to events accessible to an observer through records.
4. \textbf{Clock time}, a monotonic numerical parametrization of that reconstructed order, restricted to a clock's history.

We have distinguished three questions that are often conflated: what orders events, how is that order represented, and how is it parametrized. We have shown that the framework is consistent with relativity and acknowledges its debt to causal set theory.

\textbf{Thesis:}
\[
\boxed{\text{A numerical temporal parameter need not be fundamental}}
\]
\[
\boxed{\text{if an appropriate relational order can be constructed from physical events.}}
\]

The construction presented here is not a theory of time in the dynamical sense. It is a proof of concept that the functional role of temporal ordering can be reconstructed from relational physical structure without introducing a primitive numerical time parameter. Whether such a construction can be embedded into a complete relativistic quantum theory remains an open question.

I started this paper by saying that fifty years ago, I stood at a lectern, looked at a \( t \), and did not ask the question.

This paper is the question I finally asked.

The question is now in the open.

\section*{A Note on the Origin of This Paper}

This paper began with a question I was too afraid to ask for fifty years.

The writing of it was never a passage from confusion to clarity. It was a slow, difficult effort to lay that confusion out honestly, in full view, so that it could be read, examined, and judged.

By the end, I found myself no closer to knowing what time is. But I had reached a place I had never dared to stop before---the boundary where I could finally say: I do not know, and that is where the question belongs.

Perhaps some questions are not meant to be answered. They are meant to be carried, in good faith, to the edge of what we can say---and left there for others to continue.

This paper is a record of that carrying. It does not claim to have solved time. It only claims to have put the question on paper, cleanly and without evasion.

What happens next belongs to those who come after.

\section*{Acknowledgements}

The author thanks his students and colleagues for tolerating a teacher who sometimes stopped in the middle of a lecture and stared at a variable on the blackboard. He thanks his family for patience, his Labrador Retriever for demonstrating that time is measured in walks, and the developers of DeepSeek for engaging in a dialogue that helped clarify what had been confusion for fifty years. He also thanks the unnamed auntie from his childhood, who used time perfectly well without ever defining it---and thereby taught him something that physics textbooks did not.

\begin{thebibliography}{99}

\bibitem{bombelli1987}
Bombelli, L., Lee, J., Meyer, D., \& Sorkin, R. D. (1987). Space-time as a causal set. \textit{Physical Review Letters}, 59(5), 521.

\bibitem{surya2019}
Surya, S. (2019). The causal set approach to quantum gravity. \textit{Living Reviews in Relativity}, 22(1), 5.

\bibitem{einstein1905}
Einstein, A. (1905). Zur Elektrodynamik bewegter K\"orper. \textit{Annalen der Physik}, 322(10), 891-921.

\bibitem{einstein1916}
Einstein, A. (1916). Die Grundlage der allgemeinen Relativit\"atstheorie. \textit{Annalen der Physik}, 354(7), 769-822.

\bibitem{maxwell1865}
Maxwell, J. C. (1865). A dynamical theory of the electromagnetic field. \textit{Philosophical Transactions of the Royal Society of London}, 155, 459-512.

\bibitem{newton1687}
Newton, I. (1687). \textit{Philosophi\ae Naturalis Principia Mathematica.}

\bibitem{schrodinger1926}
Schr\"odinger, E. (1926). Quantisierung als Eigenwertproblem. \textit{Annalen der Physik}, 384(4), 361-376.

\bibitem{pagewootters1983}
Page, D. N., \& Wootters, W. K. (1983). Evolution without evolution: Dynamics described by stationary observables. \textit{Physical Review D}, 27(12), 2885.

\bibitem{rovelli1993}
Rovelli, C. (1993). Statistical mechanics of gravity and the thermodynamical origin of time. \textit{Classical and Quantum Gravity}, 10(8), 1567.

\bibitem{kuchar1992time}
Kucha\v{r}, K. V. (1992). Time and interpretations of quantum gravity. In \textit{Proceedings of the 4th Canadian Conference on General Relativity and Relativistic Astrophysics} (pp. 211-314). World Scientific.

\bibitem{isham1993canonical}
Isham, C. J. (1993). Canonical quantum gravity and the problem of time. In \textit{Integrable Systems, Quantum Groups, and Quantum Field Theories} (pp. 157-287). Springer.

\bibitem{kiefer2012}
Kiefer, C. (2012). \textit{Quantum Gravity} (3rd ed.). Oxford University Press.

\bibitem{barbour1999}
Barbour, J. (1999). \textit{The End of Time: The Next Revolution in Physics.} Oxford University Press.

\bibitem{rovelli2018}
Rovelli, C. (2018). \textit{The Order of Time.} Riverhead Books.

\bibitem{smolin2013}
Smolin, L. (2013). \textit{Time Reborn.} Houghton Mifflin Harcourt.

\end{thebibliography}

\end{document}